% ==============================================================================
% 254_Zwei_Ordnungsprinzipien_En_ch.tex
% Kapitelinhalt – einzubinden via Wrapper
% Autor: Johann Pascher · Mai 2026
% ==============================================================================

\section*{Abstract}

In the physics of black holes, two ordering principles appear that seem to contradict each other: the \textbf{resonance principle} (nature selects resonant configurations, ξ-recursion) and the \textbf{thermodynamic principle} (entropy seeks its maximum, equilibration). This document shows that the apparent contradiction is resolved through a separation of levels --- the two principles operate on different layers --- and that at the evaporation endpoint they are not merely compatible but mutually enforcing. A central geometric theorem (area-bound argument) is derived. The connection to Vopson's Second Law of Infodynamics is characterised as an honest translation --- same phenomenology, different mechanisms. Three open proof gaps are explicitly identified.

\tableofcontents
\newpage

% ------------------------------------------------------------------------------
\section{The Two Ordering Principles}

\subsection{The resonance principle}

In FFGFT, the ξ-recursion on the 4D identification torus $T^4$ selects the admissible winding configurations. Only configurations satisfying the recursion condition are stable --- particles have discrete masses, not a continuum. The ξ-principle is order-creating: it generates structure, selects the resonant path, and defines the admissible state space before dynamics operate (cf.\ Dok.~201, Dok.~250).

\subsection{The thermodynamic principle}

Thermodynamics seeks maximum entropy: disorder increases, structures equilibrate, everything flows toward the averaged state. This appears directly opposed to the resonance principle.

\begin{wichtig}[The apparent contradiction]
The resonance principle favours order and structure. The thermodynamic principle favours equilibration and disorder. Both appear to act simultaneously in nature. How is this possible?
\end{wichtig}

% ------------------------------------------------------------------------------
\section{The Area-Bound Theorem}

\subsection{The argument chain}

The following argument is a \textbf{theorem}, not a postulate. It follows from three established quantities:

\begin{enumerate}
  \item \textbf{Holographic bound (Bekenstein):} The maximum information a surface can carry is proportional to its area: $N_\text{max} \propto A$.
  \item \textbf{Hawking evaporation:} $A \propto M^2$ shrinks with mass.
  \item \textbf{FFGFT lower bound:} $L_0 = \xipar \cdot \ell_P$ sets a minimum length, hence a minimum area --- the state space ends at $N = 0$.
\end{enumerate}

\begin{keyresult}[Area-bound theorem]
From (1), (2), and (3) it follows necessarily: the \textbf{possible information capacity} of the horizon surface must decrease during evaporation. This is not postulated; it follows from the geometry alone.
\begin{equation}
  N_\text{max}(M) \propto M^2 \quad \Rightarrow \quad \frac{dN_\text{max}}{dt} < 0 \quad \text{(while the BH evaporates)}
\end{equation}
\end{keyresult}

\subsection{Concentration as a consequence}

Adding topological charge conservation $Q$ to the theorem:

\begin{erkenntnis}[Forced concentration]
\begin{itemize}
  \item The \emph{possible} information (capacity $N_\text{max}$) decreases --- geometrically forced.
  \item The \emph{actual} topological charge $Q$ is preserved --- topological invariance.
  \item Therefore $Q$ must be packed into ever fewer degrees of freedom: \textbf{concentration is enforced}.
\end{itemize}
At the remnant ($N = 0$): maximum charge concentration at minimum capacity. This is the resonance principle --- not chosen, but inevitable.
\end{erkenntnis}

\subsection{Entropy and resonance as two sides of one process}

The thermodynamic principle drives evaporation (radiation entropy increases, $dS_\text{rad}/dt > 0$). In doing so the area shrinks. The area-bound theorem then forces the concentration. Entropy and the resonance principle are therefore not opponents:

\begin{erkenntnis}[No contradiction, but a causal chain]
The entropic principle generates the conditions under which the resonance principle becomes inevitable. The one causes the other.
\end{erkenntnis}

% ------------------------------------------------------------------------------
\section{Python-Verified Key Numbers}

\subsection{Counter-directional behaviour: N\_node vs.\ radiation entropy}

Numerical verification (script \texttt{infodynamik.py}):

\begin{equation}
  N_\text{node} \propto M^2 \quad \Rightarrow \quad \frac{dN_\text{node}}{dt} < 0 \quad \text{always}
\end{equation}
\begin{equation}
  \frac{dS_\text{rad}}{dt} > 0 \quad \text{(Page factor 4/3: radiation carries away more entropy than the BH loses)}
\end{equation}

The counter-directional behaviour is exact and computable. For $M = M_\odot$:
\begin{equation}
  \frac{k_B T_H}{E_\text{max}} \approx 5.8 \times 10^{-44}, \qquad N_\text{node}(M_\odot) \approx 1.5 \times 10^{77}
\end{equation}

At the remnant ($M_\text{crit} \approx 1.15 \times 10^{-13}$ kg): $N_\text{node} \approx 5 \times 10^{-10}$ --- capacity is effectively zero, charge maximally concentrated.

\subsection{Scaling comparison: capacity vs.\ naive charge count}

Testing the consistency condition $Q \leq N_\text{max}$ with the naive count $Q = M/m_e$ (script \texttt{kapazitaet.py}):

\begin{center}
\begin{tabular}{lll}
\toprule
$M$ & $N_\text{max}/Q$ & Status \\
\midrule
$1\,M_\odot$ & $6.9 \times 10^{16}$ & consistent \\
$10^{20}$ kg & $3.5 \times 10^{6}$ & consistent \\
$10^{11}$ kg & $3.5 \times 10^{-3}$ & scissors cross \\
\bottomrule
\end{tabular}
\end{center}

\begin{wichtig}[Meaning of the crossing]
The naive count $Q = M/m_e$ exceeds $N_\text{max}$ at $M \approx 2.9 \times 10^{13}$ kg --- this is \emph{not} a physical contradiction but the proof that $M/m_e$ is the \textbf{wrong account}. It counts infallen particles, not topological invariants. The accounts are not the same.
\end{wichtig}

This yields a testable constraint: the correct topological charge $Q$ must grow at most as $M^2$ (like the capacity) for $Q \leq N_\text{max}$ to hold throughout. This is a restriction on FFGFT, not a free choice.

% ------------------------------------------------------------------------------
\section{Connection to Vopson's Infodynamics}

\subsection{Vopson's Second Law}

Vopson's Second Law of Infodynamics states: the information entropy of a system decreases or remains constant while thermodynamic entropy increases. The two entropies run in opposite directions.

\subsection{FFGFT reproduces the phenomenology}

In FFGFT quantities:
\begin{itemize}
  \item Decreasing quantity: $N_\text{max}(M) \propto M^2$ --- capacity falls.
  \item Increasing quantity: $S_\text{rad}$ --- thermodynamic entropy of the radiation rises.
\end{itemize}

The directions are the same as in Vopson's law. \textbf{But the mechanism differs}:

\begin{center}
\begin{tabular}{p{5cm}p{8cm}}
\toprule
\textbf{Vopson} & \textbf{FFGFT} \\
\midrule
Information entropy decreases as a \emph{postulated law of nature} & Capacity $N_\text{max}$ decreases as a \emph{geometric theorem} (area + Bekenstein + $L_0$) \\
Mechanism: information has mass, drives minimisation & Mechanism: area bound + conservation + lower bound \\
Counts: information-bearing particle states & Counts: surface nodes $\propto A \propto M^2$ \\
\bottomrule
\end{tabular}
\end{center}

\begin{keyresult}[Honest translation]
FFGFT reproduces the \emph{phenomenology} of Vopson's law through a different mechanism. The two accounts count different things and converge only in behaviour. The connection is a \textbf{translation}, not an identity.

This is a stronger epistemic position than Vopson's: FFGFT postulates nothing but derives the counter-directional behaviour from the geometry.
\end{keyresult}

% ------------------------------------------------------------------------------
\section{Three Open Proof Gaps}

These are named explicitly, because claim hygiene requires the obligation to self-limit.

\subsection{Gap 1: The correct topological charge}

\textbf{Open question:} What is the conserved quantity $Q$ in units comparable to $N_\text{max}$?

$Q = M/m_e$ is demonstrably wrong (exceeds the capacity). $Q$ must scale at most as $M^2$. Whether $Q$ is identical to the Bekenstein entropy itself (holographic principle in its strongest form), sub-extensive, or something else --- is the open question.

\textbf{Why this matters:} Only once $Q$ is precisely defined can the consistency condition $Q \leq N_\text{max}$ be verified as a theorem throughout evaporation --- not merely plausibilised numerically.

\subsection{Gap 2: The extremal principle}

\textbf{Open question:} Can the resonance principle be formulated as a variational statement?

The Fibonacci convergence $\varphi^{-2n}/\sqrt{5} \approx \xipar$ at $n \approx 8$ and the Lagrange stiffness term $-\frac{1}{2}m_T^2(\partial m)^2$ suggest an extremal principle. But the explicit Euler-Lagrange statement --- a proof that the ξ fixed point is a minimum of an information functional --- does not yet exist in the corpus.

\textbf{What the proof would need:} An information functional $\mathcal{I}[\Psi]$ on the configuration space of $T^4$ such that $\delta\mathcal{I}/\delta\Psi = 0$ has exactly the ξ-stable winding configurations as solutions.

\subsection{Gap 3: Definition of information entropy in FFGFT}

\textbf{Open question:} How does FFGFT define information entropy precisely?

Currently $N_\text{node}$ serves as the information measure. But $N_\text{node}$ is the holographic capacity (surface nodes), not the actually stored information. The clean definition --- what distinguishes information capacity from information entropy in FFGFT --- is still missing. This is where a precise bridge to Vopson would be possible, but is not yet fully built.

% ------------------------------------------------------------------------------
\section*{Summary}

\begin{longtable}{>{\bfseries}p{5.5cm} p{8.5cm}}
\toprule
Point & Statement \\
\midrule
Apparent contradiction & Resonance principle (order) vs.\ thermodynamic principle (equilibration) \\
Resolution & Different levels (topological vs.\ thermodynamic) \\
Theorem & Area bound + Bekenstein + $L_0$ $\Rightarrow$ capacity necessarily decreases \\
Causal chain & Entropy drives evaporation $\Rightarrow$ area shrinks $\Rightarrow$ concentration enforced \\
Vopson bridge & Same phenomenology, different mechanism: translation, not identity \\
Account difference & $Q = M/m_e$ (particles) $\neq$ $N_\text{max}$ (surface) --- two different counts \\
Open question 1 & Correct topological charge $Q$: must satisfy $Q \leq N_\text{max} \propto M^2$ \\
Open question 2 & Extremal principle: ξ fixed point as minimum of an information functional \\
Open question 3 & Clean definition of information entropy in FFGFT \\
\bottomrule
\end{longtable}
